Soit $k$ un entier naturel non nul et $S$ un ensemble fini de nombres premiers impairs. Démontrer qu'il
existe au plus une manière (à rotation et symétrie axiale près) de disposer les éléments de $S$ sur un
cercle de sorte que le produit de deux nombres placés l'un à côté de l'autre soit toujours de la forme
$x^2 + x + k$, avec $x$ entier naturel non nul.
